\section{Overview}
\label{sec:background}

This section outlines the design goals and key components of \sys.

\begin{table}[t]
    \centering
    \begin{tabular}{@{\hspace{4px}}l*{6}{@{\hspace{4px}}c}@{\hspace{4px}}}
    \toprule
     & \rot{CNR~[\citenum{checknrun}]} & \rot{QD~[\citenum{QD}]} & \rot{LC~[\citenum{lcCompression}]} & \rot{DDNN~[\citenum{deltadnn}]} & \rot{GOBO~[\citenum{gobo}]} & \rot{\sys} \\
    \midrule
    In-training & \greencheck & \greencheck & \greencheck & \greencheck & \redcross & \greencheck \\
    Low Overhead & \greencheck & \greencheck & \greencheck & \greencheck & \redcross & \greencheck \\
    Model Agnostic & \redcross & \greencheck & \greencheck & \greencheck & \greencheck & \greencheck \\
    Scalable & \greencheck & \greencheck & \greencheck & \greencheck & \redcross & \greencheck \\
    Algorithmically Transparent & \greencheck & \redcross & \greencheck & \greencheck & \greencheck & \greencheck \\
    \midrule
    Dynamic Configuration Management & \redcross & \redcross & \redcross & \greencheck & \redcross & \greencheck \\
    Non-uniform Quantization & \redcross & \redcross & \greencheck & \redcross & \greencheck & \greencheck \\
    Quantization-aware Delta Encoding & \redcross & \greencheck & \greencheck & \greencheck & \redcross & \greencheck \\
    \bottomrule
    \end{tabular}
    \vspace{0.5em}
    \caption{
        \small Comparison of checkpoint compression systems.
    }
    \label{tab:system-featurs}
    \end{table}
    


\minihead{Design Goals}
\sys aims to preserve model accuracy under multiple restores from compressed checkpoints while minimizing storage overhead. Additional design goals include:

\begin{itemize}
\item Low Overhead: Compression should impose minimal overhead ($<0.5$\%) on training runtime.
\item Model Agnostic: The quantization algorithm should handle diverse model architectures without requiring model-specific adjustments.
\item Scalable: Ability to process models with hundreds of millions to billions of parameters efficiently.
\item Algorithmically Transparent: No interference with the original training algorithm, unlike quantization-aware training methods.
\end{itemize}

\autoref{tab:system-featurs} compares the main features of \sys with existing checkpoint compression systems. A detailed discussion is provided in \S~\ref{sec:related}.

\minihead{System Overview}
\sys processes checkpoint requests in three stages, as illustrated in \autoref{fig:hld}. First, the \qcm performs a distributed search to identify an optimal quantization configuration. Next, the \quantizer compresses the model using the identified configuration. Finally, the \deltaE processes the compressed checkpoint in CPU memory. 



\noindent{\textbf{\qcm} (\S~\ref{sec:search}):} 
The quantization configuration (\autoref{fig:hld}, \autoref{tab:dq-search-space}) summarizes all quantization related parameters. Over the course of training, \sys dynamically updates the quantization configuration throughout training to maximize the compression while satisfying to user-defined  quality threshold ($\epsilon$) for maximum allowed quantization error per checkpoint. An efficient search algorithm identifies a configuration maximizing compression ratio while adhering to $\epsilon$.



\noindent{\textbf{\quantizer} (\S~\ref{sec:quantization}):} 
Given a specific quantization configuration, \sys divides all the model parameters into three categories accordingly: 
prune (set to zero if below the pruning threshold), protect (preserve in full precision if above the protection threshold), or quantize (apply non-uniform quantization).
\sys uses a novel efficient non-uniform quantization scheme using sensitivity-aware approximate K-Means clustering, where quantization granularity is determined by the number of bins used in K-Means (e.g., 8-bins=3-bits). This approach offers superior quantized models with reduced runtime overhead compared to traditional methods and allows for any number of quantization bins.

\noindent{\textbf{\deltaE} (\S~\ref{sec:delta}):} 
\sys performs lossless compression using quantization-aware delta compression, exploiting similarities between consecutive checkpoints. \deltaE employs a novel parameter rearrangement technique enabling efficient run-length encoding, reducing storage overhead by up to two orders of magnitude.
After delta-compression and run-length encoding, model parameters are encoded in a combined byte stream format.
